\IfFileExists{stacks-project.cls}{%
\documentclass{stacks-project}
}{%
\documentclass{amsart}
}

% For dealing with references we use the comment environment
\usepackage{verbatim}
\newenvironment{reference}{\comment}{\endcomment}
%\newenvironment{reference}{}{}
\newenvironment{slogan}{\comment}{\endcomment}
\newenvironment{history}{\comment}{\endcomment}

% For commutative diagrams we use Xy-pic
\usepackage[all]{xy}

% We use 2cell for 2-commutative diagrams.
\xyoption{2cell}
\UseAllTwocells

% We use multicol for the list of chapters between chapters
\usepackage{multicol}

% This is generally recommended for better output
\usepackage{lmodern}
\usepackage[T1]{fontenc}

% For cross-file-references
\usepackage{xr-hyper}

% Package for hypertext links:
\usepackage{hyperref}

% For any local file, say "hello.tex" you want to link to please
% use \externaldocument[hello-]{hello}
\externaldocument[introduction-]{introduction}
\externaldocument[conventions-]{conventions}
\externaldocument[sets-]{sets}
\externaldocument[categories-]{categories}
\externaldocument[topology-]{topology}
\externaldocument[sheaves-]{sheaves}
\externaldocument[sites-]{sites}
\externaldocument[stacks-]{stacks}
\externaldocument[fields-]{fields}
\externaldocument[algebra-]{algebra}
\externaldocument[brauer-]{brauer}
\externaldocument[homology-]{homology}
\externaldocument[derived-]{derived}
\externaldocument[simplicial-]{simplicial}
\externaldocument[more-algebra-]{more-algebra}
\externaldocument[smoothing-]{smoothing}
\externaldocument[modules-]{modules}
\externaldocument[sites-modules-]{sites-modules}
\externaldocument[injectives-]{injectives}
\externaldocument[cohomology-]{cohomology}
\externaldocument[sites-cohomology-]{sites-cohomology}
\externaldocument[dga-]{dga}
\externaldocument[dpa-]{dpa}
\externaldocument[sdga-]{sdga}
\externaldocument[hypercovering-]{hypercovering}
\externaldocument[schemes-]{schemes}
\externaldocument[constructions-]{constructions}
\externaldocument[properties-]{properties}
\externaldocument[morphisms-]{morphisms}
\externaldocument[coherent-]{coherent}
\externaldocument[divisors-]{divisors}
\externaldocument[limits-]{limits}
\externaldocument[varieties-]{varieties}
\externaldocument[topologies-]{topologies}
\externaldocument[descent-]{descent}
\externaldocument[perfect-]{perfect}
\externaldocument[more-morphisms-]{more-morphisms}
\externaldocument[flat-]{flat}
\externaldocument[groupoids-]{groupoids}
\externaldocument[more-groupoids-]{more-groupoids}
\externaldocument[etale-]{etale}
\externaldocument[chow-]{chow}
\externaldocument[intersection-]{intersection}
\externaldocument[pic-]{pic}
\externaldocument[weil-]{weil}
\externaldocument[adequate-]{adequate}
\externaldocument[dualizing-]{dualizing}
\externaldocument[duality-]{duality}
\externaldocument[discriminant-]{discriminant}
\externaldocument[derham-]{derham}
\externaldocument[local-cohomology-]{local-cohomology}
\externaldocument[algebraization-]{algebraization}
\externaldocument[curves-]{curves}
\externaldocument[resolve-]{resolve}
\externaldocument[models-]{models}
\externaldocument[functors-]{functors}
\externaldocument[equiv-]{equiv}
\externaldocument[pione-]{pione}
\externaldocument[etale-cohomology-]{etale-cohomology}
\externaldocument[proetale-]{proetale}
\externaldocument[relative-cycles-]{relative-cycles}
\externaldocument[more-etale-]{more-etale}
\externaldocument[trace-]{trace}
\externaldocument[crystalline-]{crystalline}
\externaldocument[spaces-]{spaces}
\externaldocument[spaces-properties-]{spaces-properties}
\externaldocument[spaces-morphisms-]{spaces-morphisms}
\externaldocument[decent-spaces-]{decent-spaces}
\externaldocument[spaces-cohomology-]{spaces-cohomology}
\externaldocument[spaces-limits-]{spaces-limits}
\externaldocument[spaces-divisors-]{spaces-divisors}
\externaldocument[spaces-over-fields-]{spaces-over-fields}
\externaldocument[spaces-topologies-]{spaces-topologies}
\externaldocument[spaces-descent-]{spaces-descent}
\externaldocument[spaces-perfect-]{spaces-perfect}
\externaldocument[spaces-more-morphisms-]{spaces-more-morphisms}
\externaldocument[spaces-flat-]{spaces-flat}
\externaldocument[spaces-groupoids-]{spaces-groupoids}
\externaldocument[spaces-more-groupoids-]{spaces-more-groupoids}
\externaldocument[bootstrap-]{bootstrap}
\externaldocument[spaces-pushouts-]{spaces-pushouts}
\externaldocument[spaces-chow-]{spaces-chow}
\externaldocument[groupoids-quotients-]{groupoids-quotients}
\externaldocument[spaces-more-cohomology-]{spaces-more-cohomology}
\externaldocument[spaces-simplicial-]{spaces-simplicial}
\externaldocument[spaces-duality-]{spaces-duality}
\externaldocument[formal-spaces-]{formal-spaces}
\externaldocument[restricted-]{restricted}
\externaldocument[spaces-resolve-]{spaces-resolve}
\externaldocument[formal-defos-]{formal-defos}
\externaldocument[defos-]{defos}
\externaldocument[cotangent-]{cotangent}
\externaldocument[examples-defos-]{examples-defos}
\externaldocument[algebraic-]{algebraic}
\externaldocument[examples-stacks-]{examples-stacks}
\externaldocument[stacks-sheaves-]{stacks-sheaves}
\externaldocument[criteria-]{criteria}
\externaldocument[artin-]{artin}
\externaldocument[quot-]{quot}
\externaldocument[stacks-properties-]{stacks-properties}
\externaldocument[stacks-morphisms-]{stacks-morphisms}
\externaldocument[stacks-limits-]{stacks-limits}
\externaldocument[stacks-cohomology-]{stacks-cohomology}
\externaldocument[stacks-perfect-]{stacks-perfect}
\externaldocument[stacks-introduction-]{stacks-introduction}
\externaldocument[stacks-more-morphisms-]{stacks-more-morphisms}
\externaldocument[stacks-geometry-]{stacks-geometry}
\externaldocument[moduli-]{moduli}
\externaldocument[moduli-curves-]{moduli-curves}
\externaldocument[examples-]{examples}
\externaldocument[exercises-]{exercises}
\externaldocument[guide-]{guide}
\externaldocument[desirables-]{desirables}
\externaldocument[coding-]{coding}
\externaldocument[obsolete-]{obsolete}
\externaldocument[fdl-]{fdl}
\externaldocument[index-]{index}

% Theorem environments.
%
\theoremstyle{plain}
\newtheorem{theorem}[subsection]{Theorem}
\newtheorem{proposition}[subsection]{Proposition}
\newtheorem{lemma}[subsection]{Lemma}

\theoremstyle{definition}
\newtheorem{definition}[subsection]{Definition}
\newtheorem{example}[subsection]{Example}
\newtheorem{exercise}[subsection]{Exercise}
\newtheorem{situation}[subsection]{Situation}

\theoremstyle{remark}
\newtheorem{remark}[subsection]{Remark}
\newtheorem{remarks}[subsection]{Remarks}

\numberwithin{equation}{subsection}

% Macros
%
\def\lim{\mathop{\mathrm{lim}}\nolimits}
\def\colim{\mathop{\mathrm{colim}}\nolimits}
\def\Spec{\mathop{\mathrm{Spec}}}
\def\Hom{\mathop{\mathrm{Hom}}\nolimits}
\def\Ext{\mathop{\mathrm{Ext}}\nolimits}
\def\SheafHom{\mathop{\mathcal{H}\!\mathit{om}}\nolimits}
\def\SheafExt{\mathop{\mathcal{E}\!\mathit{xt}}\nolimits}
\def\Sch{\mathit{Sch}}
\def\Mor{\mathop{\mathrm{Mor}}\nolimits}
\def\Ob{\mathop{\mathrm{Ob}}\nolimits}
\def\Sh{\mathop{\mathit{Sh}}\nolimits}
\def\NL{\mathop{N\!L}\nolimits}
\def\CH{\mathop{\mathrm{CH}}\nolimits}
\def\proetale{{pro\text{-}\acute{e}tale}}
\def\etale{{\acute{e}tale}}
\def\QCoh{\mathit{QCoh}}
\def\Ker{\mathop{\mathrm{Ker}}}
\def\Im{\mathop{\mathrm{Im}}}
\def\Coker{\mathop{\mathrm{Coker}}}
\def\Coim{\mathop{\mathrm{Coim}}}

% Boxtimes
%
\DeclareMathSymbol{\boxtimes}{\mathbin}{AMSa}{"02}

%
% Macros for moduli stacks/spaces
%
\def\QCohstack{\mathcal{QC}\!\mathit{oh}}
\def\Cohstack{\mathcal{C}\!\mathit{oh}}
\def\Spacesstack{\mathcal{S}\!\mathit{paces}}
\def\Quotfunctor{\mathrm{Quot}}
\def\Hilbfunctor{\mathrm{Hilb}}
\def\Curvesstack{\mathcal{C}\!\mathit{urves}}
\def\Polarizedstack{\mathcal{P}\!\mathit{olarized}}
\def\Complexesstack{\mathcal{C}\!\mathit{omplexes}}
% \Pic is the operator that assigns to X its picard group, usage \Pic(X)
% \Picardstack_{X/B} denotes the Picard stack of X over B
% \Picardfunctor_{X/B} denotes the Picard functor of X over B
\def\Pic{\mathop{\mathrm{Pic}}\nolimits}
\def\Picardstack{\mathcal{P}\!\mathit{ic}}
\def\Picardfunctor{\mathrm{Pic}}
\def\Deformationcategory{\mathcal{D}\!\mathit{ef}}
\usepackage{cleveref}
\crefname{originallemma}{original lemma}{original lemmas}
\Crefname{originallemma}{Original lemma}{Original lemmas}
\makeatletter
\newlabel{lemma-group-scheme-addition-tangent-vectors}{{0BF5}{source}{}{sourcepointer-1}{}}
\newlabel{lemma-group-scheme-addition-tangent-vectors@cref}{{[originallemma][0][]0BF5}{[1][0][]source}{}{}{}}
\makeatother
\crefname{originallemma}{original lemma}{original lemmas}
\Crefname{originallemma}{Original lemma}{Original lemmas}
\makeatletter
\newlabel{lemma-abelian-variety-multiplication-by-d-etale}{{0BFH}{source}{}{sourcepointer-2}{}}
\newlabel{lemma-abelian-variety-multiplication-by-d-etale@cref}{{[originallemma][0][]0BFH}{[1][0][]source}{}{}{}}
\makeatother
\crefname{originaldefinition}{original definition}{original definitions}
\Crefname{originaldefinition}{Original definition}{Original definitions}
\makeatletter
\newlabel{more-algebra-definition-p-basis}{{07P1}{source}{}{sourcepointer-3}{}}
\newlabel{more-algebra-definition-p-basis@cref}{{[originaldefinition][0][]07P1}{[1][0][]source}{}{}{}}
\makeatother
\crefname{originallemma}{original lemma}{original lemmas}
\Crefname{originallemma}{Original lemma}{Original lemmas}
\makeatletter
\newlabel{more-algebra-lemma-p-basis}{{07P2}{source}{}{sourcepointer-4}{}}
\newlabel{more-algebra-lemma-p-basis@cref}{{[originallemma][0][]07P2}{[1][0][]source}{}{}{}}
\makeatother
\crefname{originallemma}{original lemma}{original lemmas}
\Crefname{originallemma}{Original lemma}{Original lemmas}
\makeatletter
\newlabel{algebra-lemma-derivative-zero-pth-power}{{031W}{source}{}{sourcepointer-5}{}}
\newlabel{algebra-lemma-derivative-zero-pth-power@cref}{{[originallemma][0][]031W}{[1][0][]source}{}{}{}}
\makeatother
\crefname{originallemma}{original lemma}{original lemmas}
\Crefname{originallemma}{Original lemma}{Original lemmas}
\makeatletter
\newlabel{lemma-degree-multiplication-by-d}{{0BFG}{source}{}{sourcepointer-6}{}}
\newlabel{lemma-degree-multiplication-by-d@cref}{{[originallemma][0][]0BFG}{[1][0][]source}{}{}{}}
\makeatother
\title{The p-torsion points of an abelian variety in characteristic p number at most p\textasciicircum{}g}
\author{The Stacks Project Authors}
\date{stacks-project:groupoids.tex:local-snapshot-20261012; GNU FDL 1.2 or later}
\begin{document}
\maketitle
\setcounter{subsection}{1}
\noindent\textit{Editorial note.} All delivered mathematical text is the author's own, copied from the pinned Stacks Project chapter file groupoids.tex; no statement or proof step was rewritten, reordered or removed. The counted span is the author environment and its own proof environment. References inside the span are handed to the assembler as navigation stubs that carry the canonical tag of the cited tag only; whether each of them is load-bearing is recorded on the stub itself.\par\medskip
\noindent\textit{Allowed standard prerequisites (editorial).}\begin{itemize}
\item Lemma lemma-group-scheme-addition-tangent-vectors (canonical tag 0BF5). cited at proof lines 2077-2079. Derivative of [p] is multiplication by p on tangent space, hence zero.
\item Definition more-algebra-definition-p-basis (07P1) and lemma more-algebra-lemma-p-basis (07P2). cited at proof lines 2089-2091. The p-basis of k(A)/k.
\item Lemma algebra-lemma-derivative-zero-pth-power (canonical tag 031W). cited at proof lines 2092-2094. Each [p]\textasciicircum{}*(t\_i) is a pth power.
\item Lemma lemma-degree-multiplication-by-d (canonical tag 0BFG). cited at proof lines 2110-2112. Degree of [p] is p\textasciicircum{}\{2g\}.
\end{itemize}\medskip
\ifdefined\hypertarget\hypertarget{sourcepointer-1}{}\fi
\noindent\textit{Source reference.} Original lemma 0BF5: sources/stacks-chow/tags.txt, source lines 14373--14373. Cited at proof lines 2077-2079 to see the derivative of [p] is multiplication by p, hence zero in characteristic p; load-bearing.\par\medskip
\ifdefined\hypertarget\hypertarget{sourcepointer-2}{}\fi
\noindent\textit{Source reference.} Original lemma 0BFH: sources/stacks-chow/tags.txt, source lines 14385--14385. Cited parenthetically at proof lines 2079-2080 as a comparison for the derivative computation; navigation for the parallel case.\par\medskip
\ifdefined\hypertarget\hypertarget{sourcepointer-3}{}\fi
\noindent\textit{Source reference.} Original definition 07P1: sources/stacks-chow/tags.txt, source lines 9781--9781. Cited at proof lines 2089-2091 to set up t\_1, ..., t\_g; load-bearing for the p-basis construction.\par\medskip
\ifdefined\hypertarget\hypertarget{sourcepointer-4}{}\fi
\noindent\textit{Source reference.} Original lemma 07P2: sources/stacks-chow/tags.txt, source lines 9782--9782. Cited at proof lines 2089-2091 alongside 07P1; load-bearing for the p-basis construction.\par\medskip
\ifdefined\hypertarget\hypertarget{sourcepointer-5}{}\fi
\noindent\textit{Source reference.} Original lemma 031W: sources/stacks-chow/tags.txt, source lines 3999--3999. Cited at proof lines 2092-2094 to see [p]\textasciicircum{}*(t\_i) has a pth root; load-bearing for the factorization.\par\medskip
\ifdefined\hypertarget\hypertarget{sourcepointer-6}{}\fi
\noindent\textit{Source reference.} Original lemma 0BFG: sources/stacks-chow/tags.txt, source lines 14384--14384. Cited at proof lines 2110-2112 to know [p] has degree p\textasciicircum{}\{2g\}, whence the degree split with pi\_2; load-bearing for the cardinality bound.\par\medskip
\begin{lemma}
\label{lemma-abelian-variety-multiplication-by-p}
Let $k$ be a field of characteristic $p > 0$. Let $A$ be an abelian variety
of dimension $g$
over $k$. The fibre of $[p] : A \to A$ over $0$ has at most
$p^g$ distinct points.
\end{lemma}

\begin{proof}
To prove this, we may and do replace $k$ by the algebraic closure.
By Lemma \ref{lemma-group-scheme-addition-tangent-vectors}
the derivative of $[p]$ is multiplication by $p$ as a map
$T_{A/k, e} \to T_{A/k, e}$ and hence is zero (compare
with proof of Lemma \ref{lemma-abelian-variety-multiplication-by-d-etale}).
Since $[p]$ commutes with translation we conclude that the derivative of $[p]$
is everywhere zero, i.e., that the induced map
$[p]^*\Omega_{A/k} \to \Omega_{A/k}$ is zero.
Looking at generic points, we find that
the corresponding map $[p]^* : k(A) \to k(A)$
of function fields induces the zero map on $\Omega_{k(A)/k}$.
Let $t_1, \ldots, t_g$ be a p-basis of $k(A)$ over $k$
(More on Algebra, Definition \ref{more-algebra-definition-p-basis} and
Lemma \ref{more-algebra-lemma-p-basis}). Then $[p]^*(t_i)$
has a $p$th root by
Algebra, Lemma \ref{algebra-lemma-derivative-zero-pth-power}.
We conclude that
$k(A)[x_1, \ldots, x_g]/(x_1^p - t_1, \ldots, x_g^p - t_g)$ is a subextension
of $[p]^* : k(A) \to k(A)$.
Thus we can find an affine open $U \subset A$ such that
$t_i \in \mathcal{O}_A(U)$ and $x_i \in \mathcal{O}_A([p]^{-1}(U))$.
We obtain a factorization
$$
[p]^{-1}(U)
\xrightarrow{\pi_1}
\Spec(\mathcal{O}(U)[x_1, \ldots, x_g]/(x_1^p - t_1, \ldots, x_g^p - t_g))
\xrightarrow{\pi_2}
U
$$
of $[p]$ over $U$. After shrinking $U$ we may assume that $\pi_1$
is finite locally free (for example by generic flatness -- actually it is
already finite locally free in our case).
By Lemma \ref{lemma-degree-multiplication-by-d} we see that
$[p]$ has degree $p^{2g}$. Since $\pi_2$
has degree $p^g$ we see that $\pi_1$ has degree $p^g$ as well.
The morphism $\pi_2$ is a universal homeomorphism hence the fibres are
singletons. We conclude that the (set theoretic) fibres of $[p]^{-1}(U) \to U$
are the fibres of $\pi_1$. Hence they
have at most $p^g$ elements. Since $[p]$ is a homomorphism of group
schemes over $k$, the fibre of $[p] : A(k) \to A(k)$ has the
same cardinality for every $a \in A(k)$ and the proof is complete.
\end{proof}

\end{document}
